\section{Through the Illusion, into the Dream}

Gornahoor has been emphasizing the ancient ideal of \emph{polis\footnote{\url{https://gornahoor.net/?p=2775\#more-2775}}} over against the modern cosmopolitanism and egalitarianism. One might note that \emph{democracy} is a code word; it certainly doesn't stand for government of the people for the people -- just ask the patriotic, conservative British Jew Glasman\footnote{\url{https://isteve.blogspot.com/2011/07/blue-labour-rapid-rise-and-faster-fall.html}} what happens when you stand with the peasant-commoners. Nor does \emph{diversity} mean what it sounds -- it in fact is code for larger ethnic voting blocks of pre-approved groups, which are strangely non-diverse. Is anyone worried that America doesn't have enough Icelanders residing on our soil? These magic words are filled with power, which aims for power, power of the basest sort. All modern political nostrums are lies; Orwellian language has triumphed. ``Fascism'' means ``something bad'' or ``not-good'', ``tolerance'' means something ``not-bad''.

Modern mindsets cannot think in any terms other than these binary ones, which usually short-circuit quickly when dealing with reality, which is dominated by the principles of co-inherence \& transcendence (to name two notables). \textbf{Unfortunately, living on the cultural capital of past times in combination with technology has allowed the illusion of success to persist, resulting in what even their own propagandists term the ``sorcerer-state\footnote{\url{http://www.washingtonpost.com/wp-srv/style/longterm/books/reviews/erosmagicandthedeath.htm}}``}.

In ancient times, the problem of \emph{the one-and-the-many}\footnote{\url{http://plato.stanford.edu/entries/problem-of-many/}} was a known mystery, about which preparation, study, and energy were spent in order to arrive at living Truth, which was then embodied in principles of order \& \emph{imperium \& transendence}. The unity of the medieval mind stemmed from these and other kinds of transcendence.

It is not chaos alone which has ruled since 1789, but lies and illusion, since there has to be a world order, or at least, an anti-order. In the absence of local rites and gods who are genuinely transcendent and lead to such (and here we think of the Orthodox Church's proper emphasis on local customs\footnote{\url{http://www.roadtoemmaus.net/back\_issue\_articles/RTE\_15/Brittanys\_Celtic\_Past.pdf}} as salvific, potentially), the world is governed by occult powers which are properly both Luciferic and Ahrimanic, attempting to draw the world and physical reality fully across Being into the illusion that material reality is all that is. It promises perfection and multiplicity, and will deliver the exact opposite. It, too, has a metaphysical goal, one which can be termed Anti-Christ.

The question is not whether this will end, of course it will. In wonderful words, a Greek citizen\footnote{\url{http://www.firstthings.com/blogs/spengler/2010/02/12/greece-as-political-time-bomb/}} describes what is coming:

\begin{quotex}
\emph{The common psychological traits of the corruption are what the ancients called alazoneia (brash presumption of knowledge by the ignorant) and anaischuntia (shamelessness)... One of the delusions is that there is a moral kernel in the country that we can turn to for consolation and renewal. There is no such thing. The corruption went too deep. The country is completely unprotected on the cultural and moral front. This too has not seeped in. And yet when people become desperate; when their world starts to crumble around them and all their delusions about themselves and their good life not only collapse, but do so without any legacy to fall back on and no dream to look forward to, then beware. We are in unchartered territory where Furies and Ate pilot the ship.}
\end{quotex}


The last words of the Pythian oracle were said to be ``Apollo will return a second time, this time, forever.'' The question is when and how resistance will be offered, and in what state this ``return'' will find the world. In the West, there is still a wealth of potential forms which can be marshaled.

A chief goal of traditionalists, regardless of past failures and disappointments, could properly be to separate Christianity from those who would turn it into a vehicle of this degradation. ``Decayed altars are inhabited by demons'', says Ernst Junger. Or, as Aristotle has it, ``corruptio optima pessimi''. Christianity has tremendous pitfalls and problems, and for that very reason, is a prime battleground of opportunity.

It could be argued\footnote{\url{http://www.kinism.net/index.php/anglican/more/degrees\_of\_love/}} that the sacred text and tradition of Christianity is actually the very opposite of what it has been made to be, the partner of egalitarianism:

\begin{quotex}
Recall, Andrewes' argument begins with the existence of God---and communion with Him as the very substance of neighborly love. Egalitarianism tends to reverses this order making communion with God ancillary or a secondary to universal brotherhood. Not only this, but modernism also displaces the immediate and natural obligations of man for an abstracted ``bodiless'' mankind that has no necessary ties to either person, spouse, or family. Notice the degrees posited by Andrewes are in contradistiction to egalitarianism---i.e., 1). first, the dressing and feeding of our body, 2) then our spouse, 3) then providing for our children, 4) caring for extended family or `kindreds', 5) friends and those we know, 6) compatriots, and, 7) finally, our church. Notice outward ethics, or the second tablet, really depend upon concentric relations and living these out. The reverse would otherwise make a man a louse, commanding him to love the extraordinary over the familiar. Also, we might suppose the church (\#5) is further ordered between degrees of ecumenical, province, and parish since what wreck the greater church would be if all her individual local parts were in destitute?
\end{quotex}


Evil's most effective means of stopping organic growth towards the Good is to counterfeit it, and to create confusion, such that genuine efforts are not recognized and come to nothing. This is occurring in our time and day with the Church. To continue to discredit the Church \emph{per se} is to discredit Western civilization, and thus to plunge mankind further and deeper into confusion, rendering the object of assault even weaker for the next attack.

Such a difficult position may seem unfair to those who ``wake up'', but it is nonetheless real. This is why Gornahoor has repeatedly emphasized the need for personal struggle first, the \emph{frein vital} or mastery of self-\emph{manas}. Only such a one is equipped to enter the field of struggle without ultimate danger to the self or to others. It is in this struggle, and looked at in this manner, that the purest and most ultimate ground of commonality can be found with Christianity, which begins its Beatitudes with poverty of spirit, mourning, hungering \& thirsting for righteousness, etc. Or, as Nicholas Gomez Davila\footnote{\url{https://en.wikipedia.org/wiki/Nicol\%C3\%A1s\_G\%C3\%B3mez\_D\%C3\%A1vila}} would say,

\begin{quotex}
In the Christian obsessed with ``social justice'' it isn't easy to discern whether charity is flourishing or faith is expiring.
\end{quotex}


All traditionalists and Christians should think about ways and means of personal transformation which can be (at least) potential ways and means of future union, the reconstruction of Esoteric Catholicism. As Seraphim of Sarov would say, the question is not one of \emph{largeness of faith}, but of purity of faith, or \emph{resolve}. Evola would call this \emph{will}. They are very nearly the same thing. And for those who strive, they are one.

There is an older Christianity, not yet dead, which possesses this spirit, described by Sir Walter Scott in \emph{Lay of the Last Minstrel}:

\begin{quotex}
\begin{verse}
But earthly spirit could not tell\\
The heart of them that loved so well.\\
True love's the gift which God has given\\
To man alone beneath the heaven.\\
It is not fantasy's hot fire,\\
Whose wishes, soon as granted, fly;\\
It liveth not in fierce desire,\\
With dead desire it doth not die;\\
It is the secret sympathy,\\
The silver link, the silken tie,\\
Which heart to heart, and mind to mind,\\
In body and in soul can bind.
\end{verse}
\end{quotex}


Can traditionalists not with this and with men of this heart find common cause?


\flrightit{Posted on 2011-07-29 by Logres}

